% -------------------------------------------------------------------------
% WBS ID: RES-NUM-STB
% Deliverable: Stability, Consistency and Convergence Analysis
% Status: In progress
% Evidence status: Local URANS--VOF boundedness and convergence requirements established
% Review status: Not reviewed
% Last updated: 2026-07-18
% Dependencies: RES-NUM-FRM, RES-NUM-TIM, RES-NUM-SPEC
% Downstream interfaces: RES-VER-CNV, RES-VER-MET, Software solver controls
% -------------------------------------------------------------------------

\section{RES-NUM-STB --- Stability, Consistency and Convergence Analysis}
\label{sec:res-num-stb}

\subsection{Purpose and Scope}
\label{sec:res-num-stb-purpose}

% TODO: Define what this Deliverable covers and explicitly excludes.

This Deliverable records stability, boundedness and convergence
requirements for the local three-dimensional URANS--VOF finite-volume
model. It defines numerical acceptance checks; it does not replace
physical validation against observations or experiments.

\subsection{Research Questions}
\label{sec:res-num-stb-research-questions}

% TODO: State the questions that must be answered before the Deliverable
% can be accepted.

\subsection{Terminology and Definitions}
\label{sec:res-num-stb-definitions}

% TODO: Define the technical terminology, symbols, quantities and
% classifications used in this Deliverable.

\subsection{Literature Evidence}
\label{sec:res-num-stb-literature-evidence}

% TODO: Record evidence from peer-reviewed papers, authoritative datasets,
% standards, technical reports and primary documentation.
%
% Do not create a detached literature-summary list. Explain how each source
% supports, contradicts or limits a project decision.

\textcite{hirt1981vof}, \textcite{roenby2016computational} and
\textcite{vukcevic2018sharp} are retained for VOF boundedness,
interface sharpness and finite-volume phase conservation. Source role:
`3D-VOF`. \textcite{issa1986operator} is retained for pressure
correction and residual convergence in segregated incompressible
solves. Source role: `3D-FVM`. Project conclusion: a local timestep
shall be accepted only when phase boundedness, mass conservation,
pressure correction and turbulence-variable admissibility are all
within tolerance. Limitation: exact tolerances remain validation and
implementation parameters.

\subsection{Governing Relations and Quantitative Definitions}
\label{sec:res-num-stb-governing-relations}

% TODO: Record relevant equations, dimensionless groups, empirical models,
% extraction rules or quantitative definitions.
%
% Use align environments for displayed equations.
% Every displayed equation must have a unique \label{eq:...}.
% Do not place punctuation after displayed equations.

\subsection{Discrete Formulation}
\label{sec:res-num-stb-discrete-formulation}

% TODO: Complete this RES-NUM Domain-specific subsection.

The local solver is accepted at a time level only after the segregated
VOF, momentum, pressure-correction and turbulence solves have produced
admissible fields. The discrete accepted state is:

\begin{align}
  \mathcal{Q}^{n+1}_{\mathrm{local}}
  &=
  \left\{
    \alpha^{n+1},\,
    \rho^{n+1},\,
    \mu^{n+1},\,
    \mathbf{u}^{n+1},\,
    p^{n+1},\,
    k^{n+1},\,
    \omega^{n+1},\,
    \mu_t^{n+1},\,
    \phi^{n+1}
  \right\}
  \label{eq:res-num-stb-local-accepted-state}
\end{align}

\subsection{Conservation Properties}
\label{sec:res-num-stb-conservation-properties}

% TODO: Complete this RES-NUM Domain-specific subsection.

The discrete mass residual in a cell is:

\begin{align}
  R_{m,i}
  &=
  \sum_{f\in\partial K_i}
  \phi_f
  \label{eq:res-num-stb-local-mass-residual}
\end{align}

The normalised mass residual is:

\begin{align}
  \widehat{R}_{m,i}
  &=
  \frac{
    \left|
      R_{m,i}
    \right|
  }{
    \sum_{f\in\partial K_i}
    \left|
      \phi_f
    \right|
    +
    \epsilon_{\phi}
  }
  \label{eq:res-num-stb-local-normalised-mass-residual}
\end{align}

The global phase-volume balance shall be reported using:

\begin{align}
  M_{\alpha}
  &=
  \sum_i
  V_i\alpha_i
  \label{eq:res-num-stb-local-phase-volume}
\end{align}

\subsection{Consistency and Formal Accuracy}
\label{sec:res-num-stb-consistency-formal-accuracy}

% TODO: Complete this RES-NUM Domain-specific subsection.

The local URANS--VOF branch shall report whether a result used
first-order fallback, interface clipping, non-orthogonal correction
limits, turbulence clipping or timestep rejection. A result that meets
a residual tolerance only after repeated limiting shall not be treated
as formally equivalent to an unlimited converged result.

\subsection{Stability Requirements}
\label{sec:res-num-stb-stability-requirements}

% TODO: Complete this RES-NUM Domain-specific subsection.

The VOF field shall remain bounded:

\begin{align}
  0
  \leq
  \alpha_i
  \leq
  1
  \label{eq:res-num-stb-local-alpha-bounds}
\end{align}

Mixture properties shall remain positive:

\begin{align}
  \rho_i
  &>
  0
  \label{eq:res-num-stb-local-density-positive}
  \\
  \mu_i
  &>
  0
  \label{eq:res-num-stb-local-viscosity-positive}
\end{align}

Turbulence variables shall remain admissible:

\begin{align}
  k_i
  &\geq
  0
  \label{eq:res-num-stb-local-k-positive}
  \\
  \omega_i
  &>
  0
  \label{eq:res-num-stb-local-omega-positive}
\end{align}

The pressure-correction loop shall reduce
$\widehat{R}_{m,i}$ below the configured tolerance before a timestep is
accepted.

\subsection{Mesh and Time-Step Requirements}
\label{sec:res-num-stb-mesh-time-step-requirements}

% TODO: Complete this RES-NUM Domain-specific subsection.

The timestep shall satisfy the local constraints in
\cref{eq:res-num-tim-local-timestep-minimum}. A timestep may be
rejected after the solve when any boundedness, residual or load-spike
acceptance criterion fails.

\subsection{Numerical Failure Modes}
\label{sec:res-num-stb-numerical-failure-modes}

% TODO: Complete this RES-NUM Domain-specific subsection.

Local numerical failure modes include:

\begin{itemize}
  \item loss of VOF boundedness;
  \item negative density, viscosity, $k$ or $\omega$;
  \item pressure--velocity decoupling on a collocated mesh;
  \item failure of the pressure-correction loop to reduce mass
    residuals;
  \item excessive non-orthogonal correction or wall-function violation;
  \item artificial boundary reflection contaminating the impact region;
  \item nonphysical pressure or force spikes caused by interface
    smearing, timestep size or mesh quality;
  \item inconsistent force and moment integrals on the barrier patch.
\end{itemize}

\subsection{Verification Requirements}
\label{sec:res-num-stb-verification-requirements}

% TODO: Complete this RES-NUM Domain-specific subsection.

Verification shall include timestep-refinement and pressure-corrector
studies for:

\begin{itemize}
  \item phase-volume conservation;
  \item free-surface boundedness and sharpness;
  \item cellwise and global mass residuals;
  \item pressure and velocity convergence;
  \item $k$ and $\omega$ boundedness;
  \item force and moment sensitivity to timestep and pressure
    residual tolerance.
\end{itemize}

\subsection{Candidate Approaches}
\label{sec:res-num-stb-candidate-approaches}

% TODO: Define the credible candidate approaches considered.

\subsection{Critical Comparison}
\label{sec:res-num-stb-critical-comparison}

% TODO: Compare candidates using physically and project-relevant criteria.
% Do not select an approach solely on popularity or implementation ease.

\subsection{Assumptions and Applicability}
\label{sec:res-num-stb-assumptions}

% TODO: State assumptions and the conditions under which they are valid.

\subsection{Limitations and Failure Conditions}
\label{sec:res-num-stb-limitations}

% TODO: State where the evidence, model or selected approach ceases to be
% reliable or applicable.

\subsection{Project-Specific Conclusions}
\label{sec:res-num-stb-conclusions}

% TODO: Convert the evidence into conclusions specific to the Studentship
% project. Do not merely restate literature findings.

\subsection{Selected Approach and Rationale}
\label{sec:res-num-stb-selected-approach}

% TODO: Record the selected approach only after sufficient evidence exists.
% State the alternatives rejected and the reasons for rejection.

The selected local stability policy is accept-or-reject timestep
control. A timestep is accepted only when VOF boundedness, physical
property positivity, turbulence-variable admissibility, pressure/mass
convergence and output-quality checks all pass. Otherwise the timestep
is repeated with a reduced $\Delta t$ or the simulation is marked
failed according to the configured recovery limit.

\subsection{Unresolved Decisions}
\label{sec:res-num-stb-unresolved-decisions}

% TODO: Record unresolved questions, required evidence and decision owners.

Open items are the final residual tolerances, maximum pressure
corrector count, non-orthogonal correction limit, local timestep
reduction factor, maximum rejected-timestep count and load-spike
diagnostic threshold.

\subsection{Requirements and Handoffs}
\label{sec:res-num-stb-requirements}

% TODO: State requirements passed to other Research Domains, Software,
% verification activities or later execution work.

`RES-NUM-SPEC` shall call these checks before accepting a local
timestep. `RES-VER-MET` shall distinguish numerical acceptance from
physical validation, especially for pressure, force, moment and
overtopping outputs.

\subsection{Deliverable Completion Record}
\label{sec:res-num-stb-completion-record}

% TODO: Record whether the Deliverable is:
%
% - Not started
% - In progress
% - Draft complete
% - Reviewed
% - Accepted
%
% Acceptance requires:
%
% - the research questions to be answered;
% - claims to be supported by appropriate evidence;
% - assumptions and limitations to be explicit;
% - the project-specific conclusion to be stated;
% - downstream requirements to be traceable;
% - unresolved decisions to be closed or formally deferred.
